% Part: methods
% Chapter: proofs
% Section: example-1

\documentclass[../../../include/open-logic-section]{subfiles}

\begin{document}

\olfileid[tr]{mth}{prf}{ex1}

\olsection{Bir Örnek}

İlk örneğimiz, kümelerin birleşimleri ve kesişimleri hakkındaki şu
basit olgudur. Bu örnek; tanımların açılmasını, bağlaşımların ve tümel
iddiaların kanıtlanmasını ve durumlara göre kanıtı gösterecektir.

\begin{prop}
Herhangi $A$, $B$ ve $C$ kümeleri için $A \cup (B \cap C) = (A \cup B)
\cap (A \cup C)$
\end{prop}

Haydi kanıtlayalım!{}

\begin{proof}
Herhangi $A$, $B$ ve $C$ kümeleri için $A \cup (B \cap C) = (A \cup
B) \cap (A \cup C)$ olduğunu göstermek istiyoruz.
\begin{quote}
Önce önermenin ifadesindeki ``$=$'' tanımını açarız. Kümelerin özdeş
olduğunu kanıtlamanın, onların !!{element}s bakımından aynı olduğunu
göstermek anlamına geldiğini hatırlayın. Yani $A \cup (B \cap C)$'nin
bütün !!{element}s aynı zamanda $(A \cup B) \cap (A \cup C)$'nin
!!{element}s olmalı ve bunun tersi de geçerli olmalıdır. ``Bunun
tersi'', $(A \cup B) \cap (A \cup C)$'deki her !!{element} için $A
\cup (B \cap C)$'de de bir !!a{element} bulunması gerektiği anlamına
gelir. Böylece tanımı açınca bir bağlaşımı kanıtlamak zorunda olduğumuzu
görürüz. Bunu kaydedelim:
\end{quote}
Tanım gereği $A \cup (B \cap C) = (A \cup B) \cap (A \cup C)$, ancak
ve ancak $A \cup (B \cap C)$'deki her !!{element}, $(A \cup B) \cap (A
\cup C)$'de bir !!a{element} olarak ve $(A \cup B) \cap (A \cup C)$'deki
her !!{element}, $A \cup (B \cap C)$'de bir !!a{element} olarak bulunursa geçerlidir.
\begin{quote}
Bu bir bağlaşım olduğundan her bileşeni ayrı ayrı kanıtlamalıyız.
İlkiyle başlayalım: $A \cup (B \cap C)$'deki her !!{element} için aynı
zamanda $(A \cup B) \cap (A \cup C)$'de bir !!a{element} bulunduğunu
kanıtlayalım.

Bu tümel bir iddiadır; dolayısıyla $A \cup (B \cap C)$'den keyfî bir
!!{element} ele alıp onun aynı zamanda $(A \cup B) \cap (A \cup
C)$'de bir !!a{element} olarak bulunması gerektiğini gösteririz. Bu keyfî
!!{element} için ad olarak, örneğin~$z$ değişkenini seçelim.
Kanıtımız şöyle devam eder:
\end{quote}
Önce $A \cup (B \cap C)$'deki her !!{element} için aynı zamanda $(A \cup
B) \cap (A \cup C)$'de bir !!a{element} bulunduğunu kanıtlayalım. $z
\in A \cup (B \cap C)$ olsun. $z \in (A \cup B) \cap (A \cup C)$
olduğunu göstermeliyiz.
\begin{quote}
Şimdi $\cup$ ve~$\cap$ tanımlarını açmanın zamanı. Örneğin $\cup$'ın
tanımı şudur: $A \cup B = \Setabs{z}{z \in A \text{ ya da } z \in B}$.
Tanımı ``$A \cup (B \cap C)$''ye uyguladığımızda, tanımdaki ``$B$''nin
rolünü artık ``$B \cap C$'' oynar; dolayısıyla $A \cup (B \cap C) =
\Setabs{z}{z \in A \text{ ya da } z \in B \cap C}$. Böylece $z \in A
\cup (B \cap C)$ varsayımımız şu demektir: $z \in \Setabs{z}{z \in A
\text{ ya da } z \in B \cap C}$. Ayrıca $z \in
\Setabs{z}{\dots z\dots}$ ancak ve ancak \dots $z$ \dots{} ise
geçerlidir; yani burada $z \in A$ ya da $z \in B \cap C$.
\end{quote}
$\cup$ tanımı gereği ya $z \in A$ ya da $z \in B \cap C$.
\begin{quote}
Bu bir ayrışım olduğundan durumlara göre kanıt kullanmak yararlı
olacaktır. İki durumu ele alıp her birinde hedeflediğimiz sonucun (yani
``$z \in (A \cup B) \cap (A \cup C)$'') geçerli olduğunu gösteririz.
\end{quote}
Durum 1: $z \in A$ olduğunu varsayalım.
\begin{quote}
Varsayımlarımıza dayanarak yararlanabileceğimiz pek bir şey yoktur.
Öyleyse sonuçta yararlanabileceğimiz şeye bakalım. $z \in (A \cup B)
\cap (A \cup C)$ olduğunu göstermek istiyoruz. $\cap$ tanımına göre $z
\in (A \cup B) \cap (A \cup C)$ olduğunu göstermek istiyorsak hem $(A
\cup B)$'de hem de $(A \cup C)$'de olduğunu göstermeliyiz. Oysa $z \in
A \cup B$, ancak ve ancak $z \in A$ ya da $z \in B$ ise geçerlidir ve
  (1. durumun varsayımı olarak) $z \in A$ zaten elimizdedir. Aynı akıl
  yürütmeyle---$C$'yi $B$'nin yerine koyarak---$z \in A \cup C$ sonucuna varırız.
Bu argüman ters yönde ilerledi; öyleyse akıl yürütmemizi kanıtımızın
gerektirdiği yönde kaydedelim.
\end{quote}
$z \in A$ olduğundan $z \in A$ ya da $z \in B$; dolayısıyla $\cup$
tanımı gereği $z \in A \cup B$. Benzer biçimde $z \in A \cup C$. Ancak
bu, $z \in (A \cup B) \cap (A \cup C)$ demektir; bu da $\cap$ tanımı gereğidir.
\begin{quote}
Bu, durumlara göre kanıtın ilk durumunu tamamlar. Şimdi $z \in B \cap
C$ olduğu ikinci durumda sonucu türetmek istiyoruz.
\end{quote}
Durum 2: $z \in B \cap C$ olduğunu varsayalım.
\begin{quote}
Yine iki kümenin kesişimiyle çalışıyoruz. $\cap$ tanımını uygulayalım:
\end{quote}
$z \in B \cap C$ olduğundan $z$, hem $B$'de hem de $C$'de bulunan bir
!!a{element} olmalıdır; bu, $\cap$ tanımı gereğidir.
\begin{quote}
Sonucumuza yeniden bakma zamanı. $z$'nin hem $(A \cup B)$'de hem de $(A
\cup C)$'de olduğunu göstermeliyiz. Çözüm yine doğrudan çıkar.
\end{quote}
$z \in B$ olduğundan $z \in (A \cup B)$. $z \in C$ olduğundan $z \in
(A \cup C)$ de geçerlidir. Öyleyse $z \in (A \cup B) \cap (A \cup C)$.
\begin{quote}
Burada $\cup$ ve $\cap$ tanımlarını yeniden uyguladık; ama bu tanımları
zaten hatırlattığımız ve $z$ iki kümeden birindeyse onların birleşiminde
de olduğunu zaten gösterdiğimiz için ne yaptığımızı bu kadar açık
yazmamız gerekmez.

Durumlara göre kanıtın ikinci durumunu tamamladık; artık ilk sonucumuzu
ileri sürebiliriz.
\end{quote}
Öyleyse $z \in A \cup (B \cap C)$ ise $z \in (A \cup B) \cap (A \cup C)$.
\begin{quote}
Şimdi yalnızca öteki yönü, yani $(A \cup B) \cap (A \cup C)$'deki her
!!{element} için $A \cup (B \cap C)$'de bir !!a{element} bulunduğunu
göstermek istiyoruz. Daha önce olduğu gibi, ilk kümede keyfî bir
!!{element} bulunduğunu varsayıp onun ikinci kümede olması
gerektiğini göstererek bu tümel iddiayı kanıtlarız. Ne yapmak üzere
olduğumuzu belirtelim.
\end{quote}
Şimdi $z \in (A \cup B) \cap (A \cup C)$ olduğunu varsayalım. $z \in A
\cup (B \cap C)$ olduğunu göstermek istiyoruz.
\begin{quote}
Artık $z \in (A \cup B) \cap (A \cup C)$ hipotezinden hareket ediyoruz.
Burada ilk kanıt parçasındakiyle aynı~$z$'yi kullanmamız umarız fazla
kafa karıştırıcı değildir. O parçayı bitirdiğimizde orada yaptığımız
bütün varsayımlar artık yürürlükte değildir; şimdi $z$'nin ne olduğuna
ilişkin yeni varsayımlar yapabiliriz. Bu size kafa karıştırıcı
geliyorsa, aşağıda $z$ yerine başka bir değişken koymanız yeterlidir.

$z$'nin hem $A \cup B$'de hem de $A \cup C$'de olduğunu biliyoruz; bu,
$\cap$ tanımı gereğidir. $\cup$ tanımıyla bunu daha da açabiliriz: ya $z \in
A$ ya da $z \in B$ ve ayrıca ya $z \in A$ ya da $z \in C$. Bu yine
durumlara göre bir kanıta benziyor---ancak ``ve'' durumu kafa karıştırıcı
hâle getiriyor. Bunun üç olasılık bulunduğu anlamına geldiğini
düşünebilirsiniz: $z$ ya $A$'da ya $B$'de ya da $C$'dedir. Fakat bu bir
hata olur. Dikkatli olmalıyız; öyleyse her ayrışımı sırayla ele alalım.
\end{quote}
$\cap$ tanımı gereği $z \in A \cup B$ ve $z \in A \cup C$. $\cup$
tanımı gereği $z \in A$ ya da $z \in B$. Durumları ayıralım.
\begin{quote}
İlk ayrışıma odaklandığımız için, $A \cup C$'yi açarak elde ettiğimiz
ikinci ayrışımı henüz kullanmadık. Aslında ona henüz ihtiyacımız yok.
İlk durum $z \in A$'dır ve bir kümede bulunan bir !!a{element}, o kümenin
herhangi başka bir kümeyle birleşiminde de bir !!a{element} olarak bulunur. Dolayısıyla 1.
durum kolaydır:
\end{quote}
Durum 1: $z \in A$ olduğunu varsayalım. Buradan $z \in A \cup (B \cap
C)$ çıkar.
\begin{quote}
Şimdi ikinci duruma, $z \in B$'ye geçelim. Burada ikinci $\cup$'ı açıp
bir kez daha durumlara göre kanıt yapacağız:
\end{quote}
Durum 2: $z \in B$ olduğunu varsayalım. $z \in A \cup C$ olduğundan ya
$z \in A$ ya da $z \in C$. Durumları daha da ayıralım:

Durum 2a: $z \in A$. O zaman yine $z \in A \cup (B \cap C)$.
\begin{quote}
Pekâlâ, bu biraz tuhaftı. Bu durum için~$z \in B$ varsayımına gerçekte
ihtiyacımız olmadı; ama bunda bir sakınca yoktur.
\end{quote}
Durum 2b: $z \in C$. O zaman $z \in B$ ve $z \in C$; dolayısıyla $z
\in B \cap C$ ve sonuç olarak $z \in A \cup (B \cap C)$.
\begin{quote}
Bu, durumlara göre iki kanıtı da sonuçlandırır; böylece ikinci yarıyı
bitirdik.
\end{quote}
Öyleyse $z \in (A \cup B) \cap (A \cup C)$ ise $z \in A \cup (B \cap C)$.
\end{proof}

\end{document}
